\section{A periodic coefficient through the prime target}
\label{sec:prime-target-transfer}

Fix a nonzero integer $h$ and a function $W\in C_c^1((0,\infty))$.
Write
\begin{equation}\label{eq:I-W}
 \cI_W=\int_0^\infty W(t)\,dt.
\end{equation}
All constants below may depend on $h$ and $W$.  The parameter $h$ is fixed;
in particular, $h=2$ is permitted.  The assertion is nevertheless a theorem
about a periodic coefficient in a prime-target form, not a prime-pair
theorem.

Let $x=MN$, where $M,N\geq2$.  For an odd integer $q$ coprime to $h$, let
$\beta:G_q\to\C$, extend it periodically to the integers, and set it equal
to zero on the nonunit residue classes.  For $(m,hq)=1$, define the row
\begin{equation}\label{eq:prime-target-row}
 \cR_{m,h}^{(N)}(\beta)
 =\sum_{n\geq1}\beta(n)\Lambda(mn+h)W(n/N).
\end{equation}

\subsection{The rowwise transfer theorem}

We first state the classical input in the form used in the proof.  Put
\begin{equation}\label{eq:Delta-definition}
 \Delta(Y;r)=
 \max_{\substack{a\bmod r\\(a,r)=1}}
 \max_{y\leq Y}
 \left|
  \sum_{\substack{n\leq y\\n\equiv a\pmod r}}\Lambda(n)
  -\frac{y}{\varphi(r)}
 \right|.
\end{equation}
For every $L>0$, the maximal Bombieri--Vinogradov theorem gives a
$B_0=B_0(L)>0$ such that
\begin{equation}\label{eq:maximal-BV}
 \sum_{r\leq Y^{1/2}(\log Y)^{-B_0}}\Delta(Y;r)
 \ll_L \frac{Y}{(\log Y)^L}.
\end{equation}
This is the classical theorem
\citep{Bombieri1965,Vinogradov1965,IwaniecKowalski2004}, not a result proved
here.

\begin{theorem}[Operator-valued periodic transfer]
\label{thm:operator-transfer}
Fix $A,C>0$.  There is $B=B(A,C,h,W)>0$ such that the following holds for
all sufficiently large $x=MN$.  Suppose
\begin{equation}\label{eq:transfer-q-range}
 3\leq q\leq(\log x)^C,
 \qquad q\text{ odd},
 \qquad (q,h)=1,
\end{equation}
and
\begin{equation}\label{eq:transfer-modulus-range}
 2Mq\leq\frac{x^{1/2}}{(\log x)^B}.
\end{equation}
Uniformly for every $\beta:G_q\to\C$ with
$\|\beta\|_\infty\leq1$,
\begin{equation}\label{eq:l1-row-transfer}
 \sum_{\substack{M<m\leq2M\\(m,hq)=1}}
 \left|
  \cR_{m,h}^{(N)}(\beta)
  -\frac{mN}{\varphi(m)}\cI_W
   (\mathsf P_{q,h}\beta)(m\bmod q)
 \right|
 \ll_{A,C,h,W}\frac{x}{(\log x)^A}.
\end{equation}
Consequently, for arbitrary complex coefficients $|\alpha_m|\leq1$
supported on the same rows,
\begin{align}
 &\sum_{M<m\leq2M}\alpha_m
   \sum_{n\geq1}\beta(n)\Lambda(mn+h)W(n/N)
 \notag\\
 &\quad=
 N\cI_W
 \sum_{M<m\leq2M}\alpha_m\frac{m}{\varphi(m)}
 (\mathsf P_{q,h}\beta)(m)
 +O_{A,C,h,W}\!\left(\frac{x}{(\log x)^A}\right).
 \label{eq:bilinear-operator-transfer}
\end{align}
\end{theorem}

The theorem is stronger than a single bilinear estimate: it controls the
sum of the absolute row errors.  This is exactly the supremum of the error
in \eqref{eq:bilinear-operator-transfer} over all row coefficients of modulus
at most one.

\begin{proof}
Put
\begin{equation}\label{eq:Lambda-q}
 \Lambda_q(t)=\Lambda(t)\one_{(t,q)=1}.
\end{equation}
We first prove the result with $\Lambda_q$ in place of $\Lambda$.
Split the $n$-sum into $b\in G_q$.  If $(mb+h,q)=1$, the substitution
$t=mn+h$ gives the exact correspondence
\begin{equation}\label{eq:AP-reindexing}
 \left\{n\geq1:n\equiv b\pmod q\right\}
 \longleftrightarrow
 \left\{t\in\Z:t\equiv mb+h\pmod{mq},\ (t-h)/m\geq1\right\}.
\end{equation}
The class on the right is reduced.  Indeed,
$(mb+h,m)=(h,m)=1$ and it is coprime to $q$ by assumption.  On this class
$\Lambda_q(t)=\Lambda(t)$ whenever $t\geq1$.  Equivalently, after inserting
the transformed weight $W((t-h)/(mN))$, the target sum is a weighted prime
sum in that reduced class modulo $mq$.

Stieltjes partial summation, applied to
$W((t-h)/(mN))$, gives
\begin{align}
 &\sum_{\substack{n\geq1\\n\equiv b\pmod q}}
  \Lambda_q(mn+h)W(n/N)
 \notag\\
 &\qquad=
 \frac{mN}{\varphi(mq)}\cI_W+E_{m,b},
 \label{eq:one-residue-main}
\end{align}
where, for a constant $C_W>0$,
\begin{equation}\label{eq:one-residue-error}
 |E_{m,b}|\ll_W\Delta(C_Wx;mq).
\end{equation}
For sufficiently large $x$, the support of the transformed weight lies in
the positive $t$-axis.  The continuous main term is exactly
\[
 \frac1{\varphi(mq)}
 \int_0^\infty W((t-h)/(mN))\,dt
 =\frac{mN}{\varphi(mq)}\cI_W.
\]
If $(mb+h,q)>1$, the corresponding $\Lambda_q$ sum is identically zero.

Because $(m,q)=1$, we have
$\varphi(mq)=\varphi(m)\varphi(q)$.  Summing the main terms in
\eqref{eq:one-residue-main} against $\beta(b)$ therefore gives
\begin{align*}
 \frac{mN}{\varphi(m)\varphi(q)}\cI_W
 \sum_{\substack{b\in G_q\\(mb+h,q)=1}}\beta(b)
 =\frac{mN}{\varphi(m)}\cI_W(\mathsf P_{q,h}\beta)(m).
\end{align*}

For the errors, the triangle inequality and \eqref{eq:one-residue-error}
give
\begin{align}
 &\sum_{\substack{M<m\leq2M\\(m,hq)=1}}
  \left|
   \sum_{\substack{b\in G_q\\(mb+h,q)=1}}\beta(b)E_{m,b}
  \right|
 \notag\\
 &\qquad\ll_W
 \left(\sum_{b\in G_q}|\beta(b)|\right)
 \sum_{r\leq2Mq}\Delta(C_Wx;r).
 \label{eq:raw-BV-error}
\end{align}
Here we enlarged the moduli $mq$ to all $r\leq2Mq$.  Choose the exponent in
\eqref{eq:maximal-BV} larger than $A+C+2$ and then choose $B$ in
\eqref{eq:transfer-modulus-range} accordingly.  Since
\begin{equation}\label{eq:l1-beta-cost}
 \sum_{b\in G_q}|\beta(b)|\leq\varphi(q)\leq q\leq(\log x)^C,
\end{equation}
the right side of \eqref{eq:raw-BV-error} is
$O_{A,C,h,W}(x(\log x)^{-A})$.

It remains to restore the terms deleted in \eqref{eq:Lambda-q}.  If
$(mn+h,q)>1$ and $\Lambda(mn+h)\ne0$, then
\begin{equation}\label{eq:exceptional-prime-power}
 mn+h=p^k
\end{equation}
for some prime $p\mid q$.  On the support of the two scales, $mn\asymp_W x$.
For each $p$, only $O_W(1)$ powers $p^k$ lie in the resulting fixed-ratio
interval.  For each such power, the number of pairs $(m,n)$ is at most
$\tau(|p^k-h|)$.  Hence, for every fixed $\eta>0$, the total deleted
contribution, even after absolute values over all rows, is
\begin{equation}\label{eq:prime-power-error}
 \ll_{\eta,h,W}x^\eta\sum_{p\mid q}\log p
 \ll_{\eta,h,W}x^\eta\log(2q).
\end{equation}
The divisor bound $\tau(t)\ll_\eta t^\eta$ was used here.  Under
\eqref{eq:transfer-q-range}, \eqref{eq:prime-power-error} is absorbed by
$x(\log x)^{-A}$.  This proves \eqref{eq:l1-row-transfer}; multiplying the
row errors by $\alpha_m$ proves \eqref{eq:bilinear-operator-transfer}.
\end{proof}

\begin{remark}[The unsimplified complexity bound]
\label{rem:raw-complexity-bound}
Before imposing $q\leq(\log x)^C$, the proof gives, for every $L>0$ and in
the corresponding Bombieri--Vinogradov modulus range,
\begin{equation}\label{eq:raw-complexity-bound}
 \sum_m|\text{row error}_m|
 \ll_{L,h,W}
 \|\beta\|_{\ell^1(G_q)}\frac{x}{(\log x)^L}
 +O_{\eta,h,W}(x^\eta\log(2q)).
\end{equation}
This formula explains the polylogarithmic condition.  The classical theorem
controls the maximum of one reduced class for each modulus; summing an
arbitrary observable over all $q$ classes is not free.
\end{remark}

\begin{remark}[The combined-modulus boundary]
\label{rem:combined-modulus-boundary}
The modulus in \eqref{eq:AP-reindexing} is $mq$.  It would be incorrect to
retain the one-sided condition $M\leq x^{1/2}(\log x)^{-B}$ from
\eqref{eq:intro-one-sided} after introducing a period $q$.  The valid direct
condition is \eqref{eq:transfer-modulus-range}.  Likewise, a character of
small exact conductor inflated to a much larger ambient period still
contains the unit cutoff for that ambient period; its integer-period cost is
the ambient $q$, unless the coefficient is rewritten at its genuinely
smaller period.
\end{remark}
